\chapter{SSH keys, properly}
\label{ch:sshkeys}

\chapterbrief
  {Replace the password with a key, and make \code{ssh brewer} work from anywhere on your network.}
  {Public-key cryptography, from the user's side.}
  {A one-word login, and no passwords in your workflow.}

This is the chapter that changes how the rest of the project feels. Typing a
password every time you deploy is not merely tedious --- it is the reason people
stop deploying often, and deploying often is the entire point of
Chapter~\ref{ch:deploy}.

\section{What a key pair is}

You generate two mathematically linked files:

\begin{center}
\begin{tikzpicture}[font=\footnotesize\sffamily,
   b/.style={rounded corners=2pt, inner sep=6pt, minimum width=44mm,
             minimum height=13mm, align=center}]
  \node[b, draw=brewerred, fill=brewerred!5] (priv)
       {\textbf{private key}\\ \code{~/.ssh/id\_ed25519}\\[1pt]
        {\scriptsize\itshape never leaves your laptop}};
  \node[b, draw=brewergreen, fill=brewergreen!5, right=22mm of priv] (pub)
       {\textbf{public key}\\ \code{~/.ssh/id\_ed25519.pub}\\[1pt]
        {\scriptsize\itshape copy this anywhere}};
  \draw[{Stealth[length=4pt]}-, brewergrey, dashed] (priv) -- node[above, font=\scriptsize\itshape, color=brewergrey] {derived from} (pub);
  \node[b, draw=brewernavy, below=13mm of pub] (pi)
       {\textbf{the Pi}\\ \code{~/.ssh/authorized\_keys}};
  \draw[-{Stealth[length=4pt]}, brewergreen] (pub) -- node[right, font=\scriptsize\itshape, color=brewergrey] {copied to} (pi);
\end{tikzpicture}
\end{center}

\begin{analogy}{A padlock you hand out freely}
The \textbf{public key} is an open padlock. You can make unlimited copies and
leave them anywhere --- on the Pi, on GitHub, on a colleague's server. A padlock
is not a secret.

The \textbf{private key} is the only key that opens those padlocks. It stays in
one place and is never copied, never emailed, never pasted into a chat.

When you connect, the Pi holds up your padlock and says ``prove you can open
this''. Your laptop proves it \emph{without ever sending the key} --- and that
is the same trick, mathematically, that the DESFire card uses in Volume One,
Chapter 10. You are about to use challenge--response before you implement it.
\end{analogy}

\section{Generate the key}

\begin{shell}
# on your laptop
ssh-keygen -t ed25519 -C "martin@laptop"
\end{shell}

Three questions:

\begin{description}[leftmargin=0pt, style=nextline]
  \item[File to save the key] Press Enter for the default,
        \code{\textasciitilde/.ssh/id\_ed25519}. Only change it if you already
        have a key you want to keep.
  \item[Passphrase] Press Enter twice for none, \emph{or} set one --- see the
        box below.
  \item[Same again] Confirm.
\end{description}

\begin{trace}
Your identification has been saved in /Users/martin/.ssh/id_ed25519
Your public key has been saved in /Users/martin/.ssh/id_ed25519.pub
The key fingerprint is:
SHA256:xK3f...9dQ martin@laptop
\end{trace}

\begin{concept}{\texttt{ed25519}, and passphrases}
\textbf{Why ed25519 and not RSA?} It is a modern elliptic-curve algorithm:
shorter keys, faster, and no key-size decision to get wrong. Use RSA only if you
must talk to something ancient.

\textbf{Should you set a passphrase?} It encrypts the private key on disk, so a
stolen laptop does not hand over your servers. The cost is typing it --- but your
operating system's keychain remembers it after the first use, so in practice you
type it once per reboot.

For a home Pi with nothing valuable on it, no passphrase is a defensible choice.
For the key you will put on GitHub in Chapter~\ref{ch:git}, set one.
\end{concept}

\section{Copy the public key to the Pi}

\begin{shell}
ssh-copy-id martin@brewer.local
\end{shell}

It asks for your Pi password one last time, then:

\begin{trace}
Number of key(s) added: 1

Now try logging into the machine, with:   "ssh 'martin@brewer.local'"
\end{trace}

Test it:

\begin{shell}
ssh martin@brewer.local
\end{shell}

\textbf{No password.} That is the whole point.

\subsection*{If \texttt{ssh-copy-id} is missing (some macOS setups)}

\begin{shell}
brew install ssh-copy-id

# or do it by hand -- this is exactly what the tool does:
cat ~/.ssh/id_ed25519.pub | ssh martin@brewer.local \
  "mkdir -p ~/.ssh && chmod 700 ~/.ssh && \
   cat >> ~/.ssh/authorized_keys && chmod 600 ~/.ssh/authorized_keys"
\end{shell}

\begin{trap}{SSH refuses keys if the permissions are loose}
\code{\textasciitilde/.ssh} must be \code{700} and \code{authorized\_keys} must
be \code{600}. If they are group- or world-writable, the SSH server silently
ignores the key and falls back to asking for a password.

Silently is the operative word: no error, no log line you would notice, just a
password prompt that should not be there. If keys ``do not work'' and you have
checked everything else, check permissions:

\begin{shell}
# on the Pi
chmod 700 ~/.ssh
chmod 600 ~/.ssh/authorized_keys
\end{shell}
\end{trap}

\section{The config file: \texttt{ssh brewer}}

Typing \code{martin@brewer.local} forever is unnecessary. On your laptop, create
\code{\textasciitilde/.ssh/config}:

\begin{shell}
nano ~/.ssh/config
\end{shell}

\begin{trace}
Host brewer
    HostName brewer.local
    User martin
    IdentityFile ~/.ssh/id_ed25519

    # Keep the connection alive through Wi-Fi hiccups instead of
    # dropping you mid-deploy.
    ServerAliveInterval 30
    ServerAliveCountMax 6
\end{trace}

\begin{shell}
chmod 600 ~/.ssh/config
ssh brewer
\end{shell}

That is it. One word.

\begin{concept}{Why this is more than a shortcut}
\code{\textasciitilde/.ssh/config} is read by \emph{everything} that uses SSH ---
\code{scp}, \code{rsync}, \code{git}, and the deploy scripts in
Chapter~\ref{ch:deploy}. Define \code{brewer} once and it works everywhere:

\begin{shell}
scp file.txt brewer:~
rsync -a target/classes/ brewer:midnight-brewer/target/classes/
\end{shell}

It also means that when your Pi's address changes, you edit \emph{one line}
instead of hunting through scripts.

If mDNS is unreliable on your network, put the IP in \code{HostName} instead ---
everything else keeps working unchanged. That indirection is the real value.
\end{concept}

\section{Turn off password logins}

Now that keys work, close the door behind you.

\begin{shell}
# on the Pi
sudo nano /etc/ssh/sshd_config
\end{shell}

Find and set these three (they may exist commented out with a \code{\#} ---
remove it):

\begin{trace}
PasswordAuthentication no
PubkeyAuthentication yes
PermitRootLogin no
\end{trace}

Save (Ctrl-O, Enter), exit (Ctrl-X), then:

\begin{shell}
sudo systemctl restart ssh
\end{shell}

\begin{trap}{Keep your current session open}
\textbf{Do not close your SSH window until you have proved a new one works.}

Open a \emph{second} terminal and run \code{ssh brewer}. If it lets you in, you
are safe. If it does not, you still have the first window to undo the change.

If you close everything and locked yourself out, the recovery is physical: power
off, put the SD card in your laptop, and edit the config on the boot partition
--- or just re-flash. Not fatal, but a genuinely annoying half hour.
\end{trap}

Why bother, on a home network? Because a Pi with password login and a weak
password is exactly what automated scanners look for, and because the habit
matters more than this one machine.

\section{Useful SSH tricks}

\begin{shell}
# run one command without opening a shell
ssh brewer 'uptime'

# copy a file over
scp notes.txt brewer:~

# copy a whole folder back
scp -r brewer:~/logs ./logs

# see what SSH is actually doing, when it misbehaves
ssh -v brewer
\end{shell}

\code{ssh -v} is the diagnostic. It prints every step --- which key it offered,
what the server said --- and it is how you find out that the server ignored your
key because of file permissions.

\begin{checkpoint}{Keys work, passwords are off}
\begin{shell}
# from your laptop, in a NEW terminal
ssh brewer 'hostname && whoami'
\end{shell}

\begin{trace}
brewer
martin
\end{trace}

No password prompt, one word to connect, and the command ran remotely. You are
now set up better than most working developers.
\end{checkpoint}
